\subsection*{A complete first-zero enclosure and its archived resolution (NS-19)}

NS-19 records the strongest enclosure extracted here from the archived
comparison, and quantifies its resolution. The former signed, factor-ten
target is no longer pursued. The complete discrepancy is not assigned the
scale of a finite-compression benchmark.

\begin{lemma}[Evaluation bound from the archived Schur comparison]
\label{lem:v143-evaluator-dual}
Use the block operator and trial map of
Lemma~\ref{lem:v140-shifted-schur}. Suppose $T\succeq\delta I$,
$K\succ0$, and $V^*SV\succeq\kappa K$ with $0<\kappa<1$.
For a bounded real functional $\ell$, write $p=\ell G$ and let
$\ell_t$ be its restriction to the tail in ordinary coordinates. Put
\[
 A^2=pK^{-1}p^*,\qquad B^2=\delta^{-1}\norm{\ell_t}^2,
 \qquad
 C_\ell=\frac{A^2+B^2+2\sqrt{1-\kappa}\,AB}{\kappa}.
\]
For every even $v$ in the complete form domain,
\[
 |\ell(v)|^2\le C_\ell q_W[v].
\]
\end{lemma}
\begin{proof}
Write $v=Gx+(0,y)^t$, possible since $V$ is invertible, and put
$r=BV+TZ$, $X^2=x^*Kx$, $Y^2=q_T[y]$. The complete Schur identity
$r^*T^{-1}r=K-V^*SV$ gives
\[
 q_W[v]\ge X^2+Y^2-2\sqrt{1-\kappa}\,XY,
 \qquad |\ell(v)|\le AX+BY.
\]
The first inequality retains the cross term through its complete
inverse-weighted bound. The positive two-by-two matrix with diagonal
$1$ and off-diagonal $-\sqrt{1-\kappa}$ has determinant $\kappa$.
Applying Cauchy--Schwarz in that matrix metric gives the displayed
constant. This uses the already certified tail floor, not a new tail metric.
\end{proof}

\begin{proposition}[Complete first positive zero: an absolute bound]
\label{prop:v143-first-zero-bound}
Assume the complete nonnegative form $\mathsf W_4$ has its certified simple
real even ground, that the remainder of its even spectrum lies in
$[b,\infty)$ with $b=10^{-67}$, and that the archived Schur hypotheses
of Lemma~\ref{lem:v143-evaluator-dual} hold with the constants below.
Let $f$ be the normalized exact column $16$ of the archived $4097$-row
trial, $\rho=q_W[f]$, and $h=P_{\xi_4}f$. Assume the archived local
endpoint gates and $\widehat h'>d=0.0016$ throughout
$[\gamma_1-0.1,\gamma_1+0.1]$. All these hypotheses and the earlier-zero
exclusion below are verified at $1024$ and $1280$ bits.
Then the first positive zero $z_1(4)$ of the complete ground transform
is simple and
\begin{equation}\label{eq:v143-first-zero-absolute}
 -8.752082\cdot10^{-33}<z_1(4)-\gamma_1<8.752082\cdot10^{-33}.
\end{equation}
The interval includes zero; neither the sign of the discrepancy nor
a nonzero lower bound for its magnitude is certified.
\end{proposition}
\begin{proof}[Interval evaluation at $1024$ and $1280$ bits]
Set $L=2\log4$ and use the real functional
$\ell(v)=\int_{-L/2}^{L/2}v(x)\cos(\gamma_1x)\,dx$.
Use the same $G,K,\kappa=62629/100000$ and
$\delta=1.7940\cdot10^{-8}$ as above. The vector $p$ is evaluated
using every trial coefficient, and the solve with $K$ is certified.
For the normalized unshifted cosine basis, the head coordinates of
$\ell$ are
\[
 c_0=\frac{2\sin(\gamma_1L/2)}{\sqrt L\,\gamma_1},\qquad
 c_n=\frac{2\sqrt2\sin(\gamma_1L/2)}{\sqrt L}
       \frac{\gamma_1}{\gamma_1^2-(2\pi n/L)^2}.
\]
Parseval supplies the entire omitted evaluator tail, without truncation:
\[
 \norm{\ell_t}^2=\frac L2+\frac{\sin(\gamma_1L)}{2\gamma_1}
                    -\sum_{n=0}^{16}c_n^2.
\]
The interval calculations give
\[
 A^2\in(0.41411,0.41412),\qquad
 B^2\in(49877,49878),\qquad C_\ell\in(79920.06156,79920.06158)<79921.
\]
For the same unit trial $f$ and $h=P_{\xi_4}f$ used in
Proposition~\ref{prop:v142-ground-zero}, nonnegativity gives
$q_W[h]\le q_W[f]=\rho$. Hence
\[
 |\widehat h(\gamma_1)|<1.400334\cdot10^{-35}.
\]
The already verified $\widehat h'>0.0016$ on
$[\gamma_1-0.1,\gamma_1+0.1]$, the mean value theorem and its unique
root give the displacement budget
\[
 |z-\gamma_1|<8.752082\cdot10^{-33}<9\cdot10^{-33}.
\]

To identify this local root as the first positive zero, the verifier
also covers $[0,\gamma_1-0.1]$ with $304$ dyadic subintervals.
On each, direct sinc integrals for trial modes $0$ through $64$ give
an upper bound strictly below zero after adding both the old
complete-ground projection error and
$\sqrt L(\sum_{n=65}^{4096}|f_n|^2)^{1/2}$. This last term accounts
for every omitted coefficient of the exact finite trial; it is not
an omitted infinite eigenfunction tail. The projection estimate
already covers that entire tail. These signs exclude all earlier
nonnegative transform values. The previous local uniqueness and
simplicity gates finish the claim.

Both precisions use the same trial and inherited complete residual
bounds. A second implementation checks the two-by-two dual solve,
recomputes all earlier-zero signs, verifies exact dyadic coverage
with rational arithmetic, and checks the serialized enclosures.
The files \texttt{discrepancy\_budget\_b1024.json} and
\texttt{discrepancy\_budget\_b1280.json} under \texttt{evidence/v143/}
record the result and explicitly mark the signed target unresolved.
\end{proof}

\begin{proposition}[Resolution of the archived bounds and the missing energy direction]
\label{prop:v143-resolution}
Retain the hypotheses, normalization and functional $\ell$ of
Proposition~\ref{prop:v143-first-zero-bound}. Write
\[
 \rho=2.45361754930714\ldots\cdot10^{-75},\qquad
 C_\ell=79920.06156776\ldots,\qquad d=0.0016.
\]
The two decimal ellipses here denote the enclosures in the archived
$1024/1280$-bit reports, not exact midpoint inputs.
The direct bound and its scalar resolution budget are
\begin{equation}\label{eq:v143-direct-budget}
 |z_1(4)-\gamma_1|\le R_0:=\frac{\sqrt{C_\ell\rho}}d,
 \qquad R_0=8.75208157392882\ldots\cdot10^{-33}.
\end{equation}
In this formula the limiting energy input is the \emph{upper} budget
$q_W[h]\le\rho$, multiplied by the complete evaluator constant
$C_\ell$. It is not the ordinary projection error
$\sqrt{\rho/b}=1.5664027417\ldots\cdot10^{-4}$.
A lower bound on the ground energy alone does not reduce
\eqref{eq:v143-direct-budget}.

If an additional certified lower bound $0\le\underline\mu\le\mu_0$
is supplied, set $\Delta=\rho-\underline\mu$ and $e=f-h$. Then
\begin{align}
 \norm e^2&\le\frac{\Delta}{b-\underline\mu},\label{eq:v143-lower-norm}\\
 q_W[e]&\le\frac{b\Delta}{b-\underline\mu},\label{eq:v143-lower-energy}\\
 |\ell(f)-\ell(h)|&\le
 \sqrt{\frac{C_\ell b\Delta}{b-\underline\mu}}.
 \label{eq:v143-lower-transfer}
\end{align}
Thus allocating a root-resolution budget $\varepsilon=10^{-71}$
entirely to this transfer term requires the scalar condition
\begin{equation}\label{eq:v143-required-lower}
 \Delta\le
 \frac{(d\varepsilon)^2(b-\rho)}{C_\ell b-(d\varepsilon)^2}
 =3.20320065696758\ldots\cdot10^{-153}.
\end{equation}
The rounded condition $\rho-\underline\mu\le3.2031\cdot10^{-153}$
is sufficient for that transfer budget. This is a relative gap of
about $1.3055\cdot10^{-78}$ in the Rayleigh energy, not a lower
bound of size $10^{-153}$. No lower estimate this close to $\rho$
is present in the archive. With only the old ordinary-norm projection
inequality, the corresponding gap requirement would instead be
$9.2332480\ldots\cdot10^{-216}$.

Even the condition \eqref{eq:v143-required-lower} would not alone
put the complete zero within $10^{-71}$ of $\gamma_1$: the same
frozen trial has the certified value
\[
 1.0641\cdot10^{-40}<\widehat f(\gamma_1)<1.0642\cdot10^{-40}.
\]
For a $\gamma_1$-centered enclosure one must also control this center
value, since the resulting bound is
\[
 |z_1(4)-\gamma_1|\le
 \frac{|\widehat f(\gamma_1)|+
       \sqrt{C_\ell b\Delta/(b-\underline\mu)}}d.
\]
A Temple--Lehmann-type lower-energy estimate is therefore a missing
\emph{transfer} direction, not by itself a solution of the original
$\gamma_1$-centered discrepancy problem.
\end{proposition}
\begin{proof}
The direct bound is the preceding proposition. To make its right-hand
side at most $10^{-71}$ without changing the evaluator constant would
need an upper ground-energy budget of at most
$(d\varepsilon)^2/C_\ell=3.2032007355\ldots\cdot10^{-153}$;
a lower energy estimate does not provide this.
For the transfer alternative, spectral orthogonality gives
$\rho=\mu_0\norm h^2+q_W[e]$ and $q_W[e]\ge b\norm e^2$.
As $\norm h^2=1-\norm e^2$,
\[
 \rho\ge\underline\mu+(b-\underline\mu)\norm e^2,
 \qquad
 q_W[e]\le\Delta+\underline\mu\,q_W[e]/b.
\]
These prove the first two estimates; the evaluator lemma applied to $e$
proves the third. Squaring its value-to-zero budget and using
$b-\underline\mu=b-\rho+\Delta$ gives
\eqref{eq:v143-required-lower}. Replacing the evaluator estimate by
$|\ell(e)|\le\sqrt L\norm e$ gives the ordinary-norm threshold
$(d\varepsilon)^2(b-\rho)/(L-(d\varepsilon)^2)$.
The code \texttt{certify\_resolution.py} replays both thresholds,
the outward rounding and the frozen-trial value at both precisions.
The mean value theorem gives the last displayed enclosure. The
claimed extra lower-energy hypothesis has not been verified.
\end{proof}

\begin{proposition}[Finite-window meaning and the scope of the resolution floor]
\label{prop:v143-fixed-window-meaning}
Under Proposition~\ref{prop:v143-first-zero-bound}'s complete-operator,
tail, endpoint and derivative hypotheses, its enclosure is one
finite-window proximity statement of the type sought in CCM's
step~(b). Its scalar radius $R_0$ is the resolution floor of the
\emph{displayed archived comparison with these inputs unchanged}.
The $\lambda=4$, $N=120$ finite-compression benchmark
\[
 2.9250\cdot10^{-71}<z_{1,120}(4)-\gamma_1
                         <2.9251\cdot10^{-71}
\]
is below that resolution by more than $10^{38}$, and also below
that of the archived ordinary projection estimates. Even formally
omitting the nonnegative tail contribution from the same dual
comparison leaves the radius
\[
 \frac{\sqrt{\rho A^2/\kappa}}d
 =2.5174230434\ldots\cdot10^{-35}>10^{35}(3\cdot10^{-71}).
\]
The finite benchmark is not an enclosure of the complete discrepancy.
Consequently these statements neither assign that discrepancy a sign
or a $10^{-71}$ magnitude, nor prove that no different argument using
the same archived data could resolve it. This is a failure of the
stated bounds to resolve the benchmark scale, not a failure of the
complete-ground zero agreement.
\end{proposition}
\begin{proof}
The first enclosure has already been proved. The separate finite
benchmark was certified before the revised NS-19 task, with the
unchanged coefficient assembler and verified eigenpair/root gates;
its $1024/1280$-bit reports are retained as finite-compression evidence
under \texttt{evidence/v143/}. Comparing its outward interval to
$R_0$ and to the formally reduced head-only budget proves the
stated resolution gaps. Neither comparison is a lower bound on the
actual complete discrepancy. No optimality theorem for every
possible use of the v1.25/v1.26 data is asserted.
\end{proof}

CCM's second missing step additionally needs control of how
$k_\lambda$ approximates a scalar multiple of $\xi_\lambda$
\cite[Section 8]{ConnesConsaniMoscovici2025}. The fixed-window
enclosure supplies no cofinal rate, no convergence of further zeros,
and no locally uniform convergence of normalized transforms to
Riemann's $\Xi$. No new window or tail metric is introduced.
NS-19 closes the enclosure-and-resolution accounting; the former
signed target is not claimed. G2 and RH remain open.
